\documentclass{article}
% Details of the cover for Kindle Direct Publishing
% 1) Ink and paper type: standard color interior with white paper
% 2) Trim size: 7.5 x 9.25 in
% 3) Bleed settings: no bleed
% 4) Paperback cover finish: matte

% Trim size is 7.5x9.25 in.  Guess at 1 inch spine (worked for edition 1 which
% was 436 pages).
\usepackage[% text={16in,9.25in},
            paperwidth=16in, paperheight=9.25in,
            margin=0in,
            centering]{geometry}

\usepackage{graphicx}  % Bring in front cover graphic
\usepackage{ragged2e}  % Better ragged right for back cover text            

% This matches the font in asy/cover.asy            
\usepackage{fontspec}\setmainfont{FreeSans.otf}

% Background color for entire page
\usepackage{xcolor}
\usepackage{pagecolor}

\definecolor{backgroundcolor}{HTML}{FFD8CC}
\definecolor{underlinecolor}{HTML}{929191}
\pagecolor{backgroundcolor}

\pagestyle{empty}
\begin{document}
\thispagestyle{empty}

\setlength{\unitlength}{1in}
\noindent\begin{picture}(0,0)(1,-1)
  % Left to right it goes Back cover, spine, front cover.
  % Back Cover
  \put(1.5,-1.75){%
    \begin{picture}(6.5,0)
      \put(0,0){\fontsize{20pt}{24pt}\selectfont\textit{A liberal approach to the}}
      \put(0,-0.10){\color{underlinecolor}\rule{5.25in}{1.4pt}}
      \put(5.25,-0.47){\makebox[0in][r]{\fontsize{24pt}{28.8pt}\selectfont Theory of Computation}}
    \end{picture}
  }
  \put(1.5,-4.87){%
    \begin{minipage}{5.25in}\fontsize{15.5pt}{18.6pt}\selectfont\RaggedRight
      \setlength{\parskip}{1ex}
      This is a text for a first undergraduate Theory of Computation course.
      It covers Turing machines and the definition of computability,
      unsolvable problems including the Halting problem,
      a brief introduction to languages and grammars,
      Finite State machines,
      and computational complexity including the \textsl{P} versus
      \textsl{NP} question.
      In addition, each chapter ends with a number of brief topics.

      The approach is mathematical, with definitions and proofs.
      But the pedagogy is liberal, emphasizing naturalness and making
      connections with other experiences that students bring to the
      subject.
      This encourages them to be active learners, and to reflect on the
      results.

      There are more than eight hundred exercises, many illustrations,
      and many links to further readings.
      It is supported by worked answers to the exercises, classroom
      projection slides, and a full electronic version,
      all freely downloadable.
    \end{minipage}
  }
  % Spine
  \put(7.5,-0.75){\rotatebox{-90}{%
      \begin{picture}(0,0)
        \put(1.0,0.5){\fontsize{30pt}{36pt}\selectfont Theory of Computation}
        \put(7,0.5){\fontsize{20pt}{24pt}\selectfont Hefferon}
      \end{picture}
    }}
  % Front cover
  % \put(-0.22,-9.28){\includegraphics{asy/cover00.pdf}}
  \put(9.5,-10.2){\includegraphics{asy/cover00.pdf}}
  
\end{picture}

% No restrictions on photo: http://www.loc.gov/pictures/item/npc2007012636/
% http://www.loc.gov/pictures/resource/npcc.12637/
% Digital ID: (digital file from original) npcc 12637 http://hdl.loc.gov/loc.pnp/npcc.12637
% Reproduction Number: LC-DIG-npcc-12637 (digital file from original)

% ? https://color.adobe.com/Triadic-color-theme-9913086/edit/?copy=true
% FAQ says free to use.
% Triadic by Joseph Gruenthal:
%  Snow Drop #E7FFCC (flourescent green)
%  Oil #313030  (black)
%  Tin #929191  (grey)
%  Pale Blue #D2FCFF 
%  Pale Jasper #FFD8CC (reddish)
\end{document}

% These lines tells gnu-emacs to typeset with the given engine
% which requires Unicode encoding only (utf-8)
% ^c^t^s for toggling synctex. 
% ^-Shift-Click to move from pdf to source, Command-Shift-Click on OSX
%%% Local Variables:
%%% mode: latex
%%% TeX-engine: luatex
%%% TeX-source-correlate-method-active: synctex
%%% coding: utf-8
%%% End:
